\documentclass[12pt]{article}
\input{structure.tex}
\renewcommand{\bottomtitlespace}{3cm}

\usepackage{etoolbox}

\makeatletter
\preto{\@verbatim}{\topsep=0pt \partopsep=0pt}
\makeatother

\begin{document}

\renewcommand{\tl}{$8$ сек.}
\renewcommand{\ml}{$512$ MB}
\problem{Задача A1. Ops}
\addauthor{Автори: Емил Инджев}{2025}{06}{07}
\renewcommand{\header}{
    НАЦИОНАЛЕН ЛЕТЕН ТУРНИР ПО ИНФОРМАТИКА\\
    Пловдив, 6-8 юни 2025 г.\\
    Група A, 11-12 клас
}

На Алиса твърде често ѝ се налага да сортира списъци с $N$ числа и вече ѝ писва.
Затова си е купила машина, която да може да сортира числа вместо нея.
Машината обаче е доста лимитирана и Алиса не е сигурна как да я използва правилно.

Машината има $M=10^5$ регистъра индексирани от $0$ до $M-1$. Регистрите са $8$-битови, т.е. могат да имат стойности между $0$ и $255$.
Машината не може да въвежда или извежда данни; може само да изпълнява операции върху регистрите си.
Операциите се задават с 4 параметъра:
\begin{itemize}
\item $\mathrm{In}_1$, $\mathrm{In}_2$, $\mathrm{Out}$ и $\mathrm{Table}$:
\begin{itemize}
\item $\mathrm{In}_1$, $\mathrm{In}_2$ и $\mathrm{Out}$ са индекси на регистри (разрешено е двойка/и индекси да са равни).
\item $\mathrm{Table}$ е $256$ на $256$ таблица от числа със стойности между $0$ и $255$.
\end{itemize}
\end{itemize}
Тогава машината изпълнява следните стъпки:
\begin{itemize}
\item Машината прочита стойностите на регистрите с индекси $\mathrm{In}_1$ и $\mathrm{In}_2$: $V_1$ и $V_2$ съответно.
\item След това смята резултатът $U = \mathrm{Table}_{V_1, V_2}$.
\item Накрая записва резултатът в регистъра с индекс $\mathrm{Out}$.
\end{itemize}
С други думи, ако с $R_i$ означим стойността на регистър $i$, машината извършва операцията:
\begingroup
\setlength\abovedisplayskip{4pt}
\setlength\belowdisplayskip{0pt}
\begin{equation*}
R_{\mathrm{Out}} := \mathrm{Table}_{R_{\mathrm{In}_1}, R_{\mathrm{In}_2}}
\end{equation*}
\endgroup

Машината обаче има само една ``глава'' за опериране с регистри, а тя има много лимитиран обхват.
Затова \textbf{индексите на трите регистъра трябва да са близо един до друг}.
По-точно, разликата между всяка двойка индекси трябва да е най-много $D=26$:
\begingroup
\setlength\abovedisplayskip{4pt}
\setlength\belowdisplayskip{0pt}
\begin{equation*}
\max\left(\left\lvert \mathrm{In}_1 - \mathrm{In}_2 \right\rvert, \left\lvert \mathrm{In}_1 - \mathrm{Out} \right\rvert, \left\lvert \mathrm{In}_2 - \mathrm{Out} \right\rvert\right) \leq D=26
\end{equation*}
\endgroup

Числата на редицата на Алис са $A_0, \cdots, A_{N-1}$ и всички са до $10^{18}$. Тя забелязва, че биха стигали $60$ бита, за да запази едно число, т.е. са ѝ нужни по 8 регистъра на число ($8$ регистъра по $8$ бита са $64 \geq 60$ бита). Тъй като машината не може въвежда данни, входните данни се записват директно в паметта ѝ преди изпълнение на операциите,
а всички други регистри се запълват с произволни числа.
Подобно, тъй като тя не може и да извежда данни, накрая изходните данни се четат директно от паметта на машината,
а стойностите на всички други регистри се игнорират.

Форматът и на входните и на изходните данни е списък от числа (нека бройката им е $K$). Те се записват в/четат от първите $8K$ регистъра. Така $P$-тото от тези $K$ числа се намира в регистри $8P$ до $8P+7$ (защото всяко число заема $8$ регистъра и при входа и при изхода).
Числата са записани в $256$-ична бройна система от най-малко значима цифра към най-значима цифра.
Т.е. ако $P$-тото число е $X$, имаме $X = \sum_{i = 0}^{7} 256^i R_{8P + i}$. По-детайлно, при вход Алиса конвертира числото $X$ към $256$-ична бройна система и го записва в дадените регистри, а при изход тя чете стойностите на регистрите и ги интерпретира като цифри на число в $256$-ична бройна система (след което го конвертира към бройната система, в която тя работи).

В задачата на Алиса за сортиране на $N$ числа \textbf{в ненамаляващ ред}, входът се състои от числото $N$ последвано от $N$-те числа, които после ще трябва да бъдат сортирани: $A_0, \cdots, A_{N-1}$. Това значи, че входът ще се запише в първите $8 (N + 1)$ регистъра, а останалите регистри ще съдържат произволни стойности. Обърнете внимание, че за консистентност $N$ също ще е записано в $8$ регистъра, макар и да би могло да се събере в по-малко регистри. Например, нека $N = 3, A_0 = 25655, A_1 = 500, A_2 = 1966081$. Стойностите на първите $32$ регистъра ще са:
\begingroup
\setlength\abovedisplayskip{4pt}
\setlength\belowdisplayskip{0pt}
\begin{equation*}
3, 0, 0, 0, 0, 0, 0, 0, 55, 100, 0, 0, 0, 0, 0, 0, 244, 1, 0, 0, 0, 0, 0, 0, 1, 0, 30, 0, 0, 0, 0, 0
\end{equation*}

Изходът трябва да се състои от просто $N$-те числа $A_0, \cdots, A_{N-1}$, но сортирани в ненамаляващ ред. Алиса няма нужда да чете $N$ от изхода, защото вече знае стойността му, т.е. изходът се чете от първите $8N$ регистъра, а стойностите на останалите регистри се игнорират.
В горния пример, изходът трябва да е такъв, че стойностите на първите $24$ регистъра да са:
\begingroup
\setlength\abovedisplayskip{4pt}
\setlength\belowdisplayskip{0pt}
\begin{equation*}
244, 1, 0, 0, 0, 0, 0, 0, 55, 100, 0, 0, 0, 0, 0, 0, 1, 0, 30, 0, 0, 0, 0, 0
\end{equation*}

Вашата задача е да напишете програма, която праща такива операции към машината, че да сортира в ненамаляващ ред редицата от $N$ числа. Алиса е заето момиче и иска сортирането да се случи максимално бързо, т.е. трябва да използвате колкото можете по-малко операции.

\subsection{Детайли по имплементацията}

Трябва да имплементирате фунцията \verb|sortNumbers|:

\begin{verbatim}
void sortNumbers(int maxN)
\end{verbatim}

Тя ще бъде извикана точно веднъж.
Подава ѝ се стойност $N_\mathrm{max}$ (\verb|maxN| в кода), която е горна граница за $N$ за теста, т.е. $N \leq N_\mathrm{max}$.
Забележете, че програмата Ви не знае точната стойност на $N$ (все пак тя е записана в първите $8$ регистъра).

Вашата програма може да извиква имплементираната от журито функция \verb|applyOp|:

\begin{verbatim}
void applyOp(int input1, int input2, int output,
    const std::array<std::array<int, 256>, 256>& opTable)
\end{verbatim}

Тя изпълнява една операция както е описано по-горе.
Обърнете внимание на типа на \verb|opTable|: това е тип, който се използва както 2д масив с измерения $256$ на $256$,
но за разлика от просто \verb|int opTable[256][256]|, когато се подава към функция,
компилаторът верифицира, че измеренията съвпадат с очакваните.

Вашата програма ще се компилира заедно с грейдър предоставен от журито. Тя не трябва да имплементира \verb|main| функция или да чете от стандартния вход или да пише към стандартния изход. Трябва да включва хедър файла \verb|ops.h|. Предоставени са Ви локални грейдър и хедър файлове, които можете да използвате, за да тествате програмата си локално.

В официалната грейдинг система, програмата на журито ще изпълни Вашите операции върху $T$ подтеста и теста ще бъде преминат успешно само ако операциите Ви сортират числата във всички подтестове (подтестовете може да имат различни списъци от числа, включително и различни стойности на $N$). Обърнете внимание, че това \textbf{не} значи, че функцията Ви \verb|sortNumbers| ще бъде извикана $T$ пъти (ще бъде извикана веднъж).

\subsection{Ограничения}

\begin{itemize}
\item $1 \leq N \leq N_\mathrm{max} \leq 350$
\item $0 \leq A_i \leq 10^{18}$
\item $M = 10^5$
\item $D = 26$
\item $T \leq 30$
\end{itemize}

\subsection{Подзадачи}

\begin{table}[H]
\begin{tblr}{|Q[c,m]|Q[c,m]|Q[c,m]|Q[c,m]|}
\hline
\textbf{Подзадача} & \textbf{Точки} & $N_\mathrm{max}$ & Други ограничения \\
\hline
$1$ & $10$ & $= 2$ & $N = N_\mathrm{max}$, $A_i \leq 200$ \\
\hline
$2$ &  $9$ & $= 2$ & $N = N_\mathrm{max}$ \\
\hline
$3$ & $11$ & $\leq 150$ & $N = N_\mathrm{max}$ \\
\hline
$4$ & $11$ & $\leq 150$ & \\
\hline
$5$ & $38$ & $\leq 350$ & $N = N_\mathrm{max}$ \\
\hline
$6$ & $21$ & $\leq 350$ & \\
\hline
\end{tblr}
\caption*{Точки за дадена подзадача се получават само ако се преминат успешно всички тестове в нея и всички други съдържащи се по ограничения подзадачи.}
\end{table}

\subsection{Оценяване}

Резултатът Ви (получената от Вас част от точките) за дадена подзадача е минималният Ви резултат на някой тест предвиден за нея.
Нека $C$ е броят операции, които програмата Ви праща на машината и нека $G = 2150000$ да е целевият брой операции.
Тогава резултатът Ви $S$ за теста се определя както следва:

\begin{itemize}
\item $C \leq G$: $S = 1$
\item $C \geq 3G$: $S = 0$
\item Иначе: $S = (1 - 0.5 \times (C / G - 1))^{2.25}$
\end{itemize}

\end{document}
